\documentclass[../../main.tex]{subfiles}% 注意这里的文件路径不能用 ./main.tex ,否则用latexmk编译子文件会报错
\graphicspath{{\subfix{./image/}}{\subfix{../../image/}}} % 指定图片引用路径,后续可以直接使用图片文件名
% 如果图片存放在上级目录,则使用 ../../image/ 作为备用路径

% 例如:
% \begin{figure}[H]
% \centering
% \includegraphics[scale=0.3]{图.png}
% \caption{}
% \label{figure:图}
% \end{figure}
% 注意:上述\label{}一定要放在\caption{}之后,否则引用图片序号会只会显示??.

\begin{document}

\section{曲面上的正交标架场}

\begin{definition}
设欧氏空间$E^3$中曲面$S$的参数方程为$\boldsymbol{r}=\boldsymbol{r}(u^1,u^2)$，记
\begin{align}
\boldsymbol{r}_\alpha=\frac{\partial \boldsymbol{r}}{\partial u^\alpha},\ \alpha=1,2;\quad \boldsymbol{n}=\frac{\boldsymbol{r}_1\times\boldsymbol{r}_2}{|\boldsymbol{r}_1\times\boldsymbol{r}_2|}.
\label{eq:natural-frame-def}
\end{align}
称$\{\boldsymbol{r};\boldsymbol{r}_1,\boldsymbol{r}_2,\boldsymbol{n}\}$为曲面$S$的\textbf{自然标架场}，该标架族依赖于参数$u^1,u^2$。

若单位正交标架场$\{\boldsymbol{r};\boldsymbol{e}_1,\boldsymbol{e}_2,\boldsymbol{e}_3\}$满足$\boldsymbol{e}_1,\boldsymbol{e}_2$为曲面$S$的切向量，则称其为曲面$S$的\textbf{一阶标架场}。
\end{definition}

设曲面$S$的两个基本形式为
\begin{align}
\mathrm{I}=g_{\alpha\beta}\mathrm{d}u^\alpha\mathrm{d}u^\beta,\quad \mathrm{II}=b_{\alpha\beta}\mathrm{d}u^\alpha\mathrm{d}u^\beta.
\label{eq:two-fundamental-forms}
\end{align}
由$\mathrm{d}\boldsymbol{r}=\boldsymbol{r}_1\mathrm{d}u^1+\boldsymbol{r}_2\mathrm{d}u^2$，结合相对分量的定义，可得自然标架场的相对分量
\begin{align}
\omega^1=\mathrm{d}u^1,\quad \omega^2=\mathrm{d}u^2,\quad \omega^3=0.
\label{eq:natural-omega123}
\end{align}

由曲面论的Gauss-Weingarten公式
\begin{align*}
&\mathrm{d}\boldsymbol{r}_\alpha=\frac{\partial \boldsymbol{r}_\alpha}{\partial u^\beta}\mathrm{d}u^\beta=\Gamma^\gamma_{\alpha\beta}\mathrm{d}u^\beta\boldsymbol{r}_\gamma+b_{\alpha\beta}\mathrm{d}u^\beta\boldsymbol{n},\\
&\mathrm{d}\boldsymbol{n}=\frac{\partial \boldsymbol{n}}{\partial u^\beta}\mathrm{d}u^\beta=-b^\gamma_\beta\mathrm{d}u^\beta\boldsymbol{r}_\gamma,
\end{align*}
其中$\Gamma^\gamma_{\alpha\beta}$是度量矩阵$(g_{\alpha\beta})$的Christoffel记号，可得自然标架场其余相对分量
\begin{align}
\omega^\gamma_\alpha=\Gamma^\gamma_{\alpha\beta}\mathrm{d}u^\beta,\quad \omega^3_\alpha=b_{\alpha\beta}\mathrm{d}u^\beta,\quad \omega^\gamma_3=-b^\gamma_\beta\mathrm{d}u^\beta,\quad \omega^3_3=0,\ \alpha,\gamma=1,2.
\label{eq:natural-omega-i-j}
\end{align}

\begin{proposition}
曲面$S$自然标架场$\{\boldsymbol{r};\boldsymbol{r}_1,\boldsymbol{r}_2,\boldsymbol{n}\}$的结构方程等价于曲面的Gauss-Codazzi方程。
\end{proposition}

\begin{proof}

1. 第一组结构方程的自洽性
由$\mathrm{d}\omega^j=0$，结合$\Gamma^\gamma_{\alpha\beta},b_{\alpha\beta}$关于$\alpha,\beta$的对称性，有
\begin{align*}
\sum_{k=1}^3\omega^k\wedge\omega^\gamma_k=\sum_{\alpha,\beta=1}^2\mathrm{d}u^\alpha\wedge\Gamma^\gamma_{\alpha\beta}\mathrm{d}u^\beta=(\Gamma^\gamma_{12}-\Gamma^\gamma_{21})\mathrm{d}u^1\wedge\mathrm{d}u^2=0,\\
\sum_{k=1}^3\omega^k\wedge\omega^3_k=\sum_{\alpha,\beta=1}^2\mathrm{d}u^\alpha\wedge b_{\alpha\beta}\mathrm{d}u^\beta=(b_{12}-b_{21})\mathrm{d}u^1\wedge\mathrm{d}u^2=0.
\end{align*}
因此第一组结构方程
\begin{align*}
\mathrm{d}\omega^j=\sum_{k=1}^3\omega^k\wedge\omega^j_k,\quad 1\leqslant j\leqslant3
\end{align*}
自动成立。

2. 第二组结构方程的等价性分析
第二组结构方程为
\begin{align*}
&\mathrm{d}\omega^\gamma_\alpha=\omega^\beta_\alpha\wedge\omega^\gamma_\beta+\omega^3_\alpha\wedge\omega^\gamma_3,\quad \mathrm{d}\omega^3_\alpha=\omega^\beta_\alpha\wedge\omega^3_\beta,\\
&\mathrm{d}\omega^\gamma_3=\omega^3_3\wedge\omega^\gamma_\beta,\quad \mathrm{d}\omega^3_3=\omega^\beta_3\wedge\omega^3_\beta.
\end{align*}
对第四式，因$\mathrm{d}\omega^3_3=0$，且
\begin{align*}
\omega^\beta_3\wedge\omega^3_\beta=-b^\beta_\alpha\mathrm{d}u^\alpha\wedge b_{\beta\gamma}\mathrm{d}u^\gamma=(b_2^\beta b_{\beta1}-b_1^\beta b_{\beta2})\mathrm{d}u^1\wedge\mathrm{d}u^2=0,
\end{align*}
故第四式自然成立。
由$g_{\alpha\gamma}b^\gamma_\beta=b_{\alpha\beta}$，可得$\omega^3_\alpha=-g_{\alpha\gamma}\omega^\gamma_3$，代入可证第二、三式等价，只需验证第一、二式。

3. 第一式对应Gauss方程
直接计算得
\begin{gather*}
\mathrm{d}\omega^\gamma_\alpha=\left(\frac{\partial\Gamma^\gamma_{\alpha2}}{\partial u^1}-\frac{\partial\Gamma^\gamma_{\alpha1}}{\partial u^2}\right)\mathrm{d}u^1\wedge\mathrm{d}u^2,\\
\omega^\beta_\alpha\wedge\omega^\gamma_\beta=\left(\Gamma^\beta_{\alpha1}\Gamma^\gamma_{\beta2}-\Gamma^\beta_{\alpha2}\Gamma^\gamma_{\beta1}\right)\mathrm{d}u^1\wedge\mathrm{d}u^2,\\
\omega^3_\alpha\wedge\omega^\gamma_3=(b_{\alpha2}b_1^\gamma-b_{\alpha1}b_2^\gamma)\mathrm{d}u^1\wedge\mathrm{d}u^2.
\end{gather*}
代入结构方程第一式，整理得
\begin{align*}
\mathrm{d}\omega^\gamma_\alpha-\omega^\beta_\alpha\wedge\omega^\gamma_\beta-\omega^3_\alpha\wedge\omega^\gamma_3=-g^{\gamma\beta}\left(R_{\alpha\beta12}-b_{\alpha1}b_{\beta2}+b_{\alpha2}b_{\beta1}\right)\mathrm{d}u^1\wedge\mathrm{d}u^2,
\end{align*}
即Gauss方程
\begin{align}
R_{1212}=b_{11}b_{22}-b_{12}b_{21}.
\label{eq:gauss-equation}
\end{align}

4. 第二式对应Codazzi方程
对$\omega^3_\alpha$求外微分，得
\begin{align*}
\mathrm{d}\omega^3_\alpha=\left(\frac{\partial b_{\alpha2}}{\partial u^1}-\frac{\partial b_{\alpha1}}{\partial u^2}\right)\mathrm{d}u^1\wedge\mathrm{d}u^2,\\
\omega^\beta_\alpha\wedge\omega^3_\beta=\left(\Gamma^\beta_{\alpha1}b_{\beta2}-\Gamma^\beta_{\alpha2}b_{\beta1}\right)\mathrm{d}u^1\wedge\mathrm{d}u^2.
\end{align*}
代入结构方程第二式，得Codazzi方程
\begin{align}
\frac{\partial b_{\alpha2}}{\partial u^1}-\frac{\partial b_{\alpha1}}{\partial u^2}=\Gamma^\beta_{\alpha1}b_{\beta2}-\Gamma^\beta_{\alpha2}b_{\beta1},\quad \alpha=1,2.
\label{eq:codazzi-equation}
\end{align}
综上，自然标架场的结构方程等价于曲面的Gauss-Codazzi方程。

\end{proof}

对曲面第一基本形式$\mathrm{I}=E(\mathrm{d}u)^2+2F\mathrm{d}u\mathrm{d}v+G(\mathrm{d}v)^2$，记$g=EG-F^2$，通过Schmidt正交化可得单位正交切向量
\begin{align}
\boldsymbol{e}_1=\frac{\boldsymbol{r}_u}{\sqrt{E}},\quad \boldsymbol{e}_2=\frac{1}{\sqrt{g}}\left(-\frac{F}{\sqrt{E}}\boldsymbol{r}_u+\sqrt{E}\boldsymbol{r}_v\right),\quad \boldsymbol{e}_3=\boldsymbol{e}_1\times\boldsymbol{e}_2=\boldsymbol{n}.
\label{eq:schmidt-orth}
\end{align}
其矩阵形式为
\begin{align*}
\begin{pmatrix}\boldsymbol{e}_1\\\boldsymbol{e}_2\end{pmatrix}=\begin{pmatrix}\frac{1}{\sqrt{E}}&0\\-\frac{F}{\sqrt{Eg}}&\frac{\sqrt{E}}{\sqrt{g}}\end{pmatrix}\cdot\begin{pmatrix}\boldsymbol{r}_u\\\boldsymbol{r}_v\end{pmatrix}.
\end{align*}
由$\mathrm{d}\boldsymbol{r}=\omega^1\boldsymbol{e}_1+\omega^2\boldsymbol{e}_2$，解得一阶标架场的切向相对分量
\begin{align}
\omega^1=\sqrt{E}\mathrm{d}u+\frac{F}{\sqrt{E}}\mathrm{d}v,\quad \omega^2=\frac{\sqrt{g}}{\sqrt{E}}\mathrm{d}v,\quad \omega^3=0.
\label{eq:first-order-omega12}
\end{align}
该过程等价于对第一基本形式配平方。

\begin{theorem}
设$\omega^1,\omega^2$为依赖于$u,v$的处处线性无关的一次微分式，则存在唯一的一次微分式$\omega_2^1=-\omega_1^2$，满足
\begin{align}
\mathrm{d}\omega^1=\omega^2\wedge\omega_2^1,\quad \mathrm{d}\omega^2=\omega^1\wedge\omega_1^2.
\label{eq:structure-omega12}
\end{align}
\end{theorem}

\begin{proof}

曲面一阶标架场的相对分量为$u,v$的一次微分式，设$\omega_2^1=-\omega_1^2=p\omega^1+q\omega^2$，代入式\eqref{eq:structure-omega12}得
\begin{align*}
\mathrm{d}\omega^1=p\omega^1\wedge\omega^2,\quad \mathrm{d}\omega^2=q\omega^1\wedge\omega^2.
\end{align*}
$\mathrm{d}\omega^1,\mathrm{d}\omega^2$为二次外微分式，必为$\mathrm{d}u\wedge\mathrm{d}v$的倍数；又$\omega^1,\omega^2$线性无关，故$\omega^1\wedge\omega^2$是$\mathrm{d}u\wedge\mathrm{d}v$的非零函数倍。因此
\begin{align*}
p=\frac{\mathrm{d}\omega^1}{\omega^1\wedge\omega^2},\quad q=\frac{\mathrm{d}\omega^2}{\omega^1\wedge\omega^2},
\end{align*}
即$\omega_2^1=-\omega_1^2$由$\omega^1,\omega^2$唯一确定。

\end{proof}

由结构方程$0=\mathrm{d}\omega^3=\omega^1\wedge\omega_1^3+\omega^2\wedge\omega_2^3$，结合Cartan引理，得
\begin{align}
(\omega_1^3,\omega_2^3)=(\omega^1,\omega^2)\cdot\begin{pmatrix}a&b\\b&c\end{pmatrix}.
\label{eq:cartan-lemma-omega3}
\end{align}
第二基本形式可表示为
\begin{align}
\mathrm{II}=-\mathrm{d}\boldsymbol{r}\cdot\mathrm{d}\boldsymbol{e}_3=\omega^1\omega_1^3+\omega^2\omega_2^3=a(\omega^1)^2+2b\omega^1\omega^2+c(\omega^2)^2.
\label{eq:ii-omega-form}
\end{align}
若已知$\mathrm{II}=L(\mathrm{d}u)^2+2M\mathrm{d}u\mathrm{d}v+N(\mathrm{d}v)^2$，将$\mathrm{d}u,\mathrm{d}v$用$\omega^1,\omega^2$表示后代入，可确定待定系数$a,b,c$。

\begin{proposition}
给定曲面$S$的第一、第二基本形式，将第一基本形式配平方写为两个一次微分式$\omega^1,\omega^2$的平方和，则$\omega^1,\omega^2,\omega^3=0$为某一阶标架场的相对分量；$\omega_2^1=-\omega_1^2$由定理唯一确定，$\omega_1^3=-\omega_3^1,\omega_2^3=-\omega_3^2$由第二基本形式唯一确定，剩余结构方程等价于Gauss-Codazzi方程。
\end{proposition}

\begin{example}
将曲面$S$的第一基本形式$\mathrm{I}=E(\mathrm{d}u)^2+2F\mathrm{d}u\mathrm{d}v+G(\mathrm{d}v)^2$配平方为
\begin{align*}
\mathrm{I}=\left(\frac{\sqrt{g}}{\sqrt{G}}\mathrm{d}u\right)^2+\left(\frac{F}{\sqrt{G}}\mathrm{d}u+\sqrt{G}\mathrm{d}v\right)^2,
\end{align*}
求其对应的一阶标架场。
\end{example}

\begin{solution}

命
\begin{align*}
\omega^1=\frac{\sqrt{g}}{\sqrt{G}}\mathrm{d}u,\quad \omega^2=\frac{F}{\sqrt{G}}\mathrm{d}u+\sqrt{G}\mathrm{d}v,
\end{align*}
即
\begin{align*}
(\omega^1,\omega^2)=(\mathrm{d}u,\mathrm{d}v)\cdot\begin{pmatrix}\frac{\sqrt{g}}{\sqrt{G}}&\frac{F}{\sqrt{G}}\\0&\sqrt{G}\end{pmatrix}.
\end{align*}
由$\mathrm{d}\boldsymbol{r}=(\omega^1,\omega^2)\cdot\begin{pmatrix}\boldsymbol{e}_1\\\boldsymbol{e}_2\end{pmatrix}=(\mathrm{d}u,\mathrm{d}v)\cdot\begin{pmatrix}\boldsymbol{r}_u\\\boldsymbol{r}_v\end{pmatrix}$，得
\begin{align*}
\begin{pmatrix}\boldsymbol{r}_u\\\boldsymbol{r}_v\end{pmatrix}=\begin{pmatrix}\frac{\sqrt{g}}{\sqrt{G}}&\frac{F}{\sqrt{G}}\\0&\sqrt{G}\end{pmatrix}\cdot\begin{pmatrix}\boldsymbol{e}_1\\\boldsymbol{e}_2\end{pmatrix},
\end{align*}
即
\begin{align*}
\begin{pmatrix}\boldsymbol{e}_1\\\boldsymbol{e}_2\end{pmatrix}=\begin{pmatrix}\frac{\sqrt{G}}{\sqrt{g}}&-\frac{F}{\sqrt{Gg}}\\0&\frac{1}{\sqrt{G}}\end{pmatrix}\cdot\begin{pmatrix}\boldsymbol{r}_u\\\boldsymbol{r}_v\end{pmatrix}.
\end{align*}

\end{solution}

\begin{example}
设$(u,v)$为曲面$S$的正交参数系，$\mathrm{I}=E(\mathrm{d}u)^2+G(\mathrm{d}v)^2$，$\mathrm{II}=L(\mathrm{d}u)^2+2M\mathrm{d}u\mathrm{d}v+N(\mathrm{d}v)^2$，求曲面$S$上一个一阶标架场及其相对分量。
\end{example}

\begin{solution}

由$\mathrm{I}=(\sqrt{E}\mathrm{d}u)^2+(\sqrt{G}\mathrm{d}v)^2$，取
\begin{align*}
\omega^1=\sqrt{E}\mathrm{d}u,\quad \omega^2=\sqrt{G}\mathrm{d}v,\quad \omega^3=0,
\end{align*}
对应$\boldsymbol{e}_1=\frac{1}{\sqrt{E}}\boldsymbol{r}_u,\ \boldsymbol{e}_2=\frac{1}{\sqrt{G}}\boldsymbol{r}_v$。

设$\omega_2^1=-\omega_1^2=p\mathrm{d}u+q\mathrm{d}v$，计算得
\begin{align*}
\mathrm{d}\omega^1=-(\sqrt{E})_v\mathrm{d}u\wedge\mathrm{d}v=\omega_2^1\wedge\omega^2=p\sqrt{G}\mathrm{d}u\wedge\mathrm{d}v,\\
\mathrm{d}\omega^2=(\sqrt{G})_u\mathrm{d}u\wedge\mathrm{d}v=\omega^1\wedge\omega_1^2=q\sqrt{E}\mathrm{d}u\wedge\mathrm{d}v,
\end{align*}
故
\begin{align*}
p=-\frac{(\sqrt{E})_v}{\sqrt{G}},\quad q=\frac{(\sqrt{G})_u}{\sqrt{E}},\quad \omega_2^1=-\frac{(\sqrt{E})_v}{\sqrt{G}}\mathrm{d}u+\frac{(\sqrt{G})_u}{\sqrt{E}}\mathrm{d}v.
\end{align*}

设$\omega_1^3=-\omega_3^1=a\omega^1+b\omega^2,\ \omega_2^3=-\omega_3^2=b\omega^1+c\omega^2$，代入第二基本形式得
\begin{align*}
\mathrm{II}=aE(\mathrm{d}u)^2+2b\sqrt{EG}\mathrm{d}u\mathrm{d}v+cG(\mathrm{d}v)^2,
\end{align*}
对比得$L=aE,\ M=b\sqrt{EG},\ N=cG$，即
\begin{align*}
a=\frac{L}{E},\quad b=\frac{M}{\sqrt{EG}},\quad c=\frac{N}{G},
\end{align*}
因此
\begin{align*}
\omega_1^3=\frac{1}{\sqrt{E}}(L\mathrm{d}u+M\mathrm{d}v),\quad \omega_2^3=\frac{1}{\sqrt{G}}(M\mathrm{d}u+N\mathrm{d}v).
\end{align*}

\end{solution}

\begin{example}
已知曲面$S$的第一基本形式$\mathrm{I}$和第二基本形式$\mathrm{II}$分别为
\begin{align*}
\mathrm{I} &= (1+u^2)(\mathrm{d}u)^2 + u^2(\mathrm{d}v)^2, \\
\mathrm{II} &= \frac{1}{\sqrt{1+u^2}}(\mathrm{d}u)^2 + \frac{u^2}{\sqrt{1+u^2}}(\mathrm{d}v)^2.
\end{align*}
求曲面$S$上的一个一阶标架场的相对分量，并验证Gauss-Codazzi方程成立。
\begin{solution}

将第一基本形式配平方和：
\begin{align*}
\mathrm{I} = \big(\sqrt{1+u^2}\mathrm{d}u\big)^2 + (u\mathrm{d}v)^2,
\end{align*}
取
$$\omega^1 = \sqrt{1+u^2}\mathrm{d}u,\quad \omega^2 = u\mathrm{d}v,\quad \omega^3=0.$$
求外微分：
\begin{align*}
\mathrm{d}\omega^1 = 0, \quad
\mathrm{d}\omega^2 = \mathrm{d}u\land\mathrm{d}v, \quad
\omega^1\land\omega^2 = u\sqrt{1+u^2}\mathrm{d}u\land\mathrm{d}v.
\end{align*}
设联络形式满足
\begin{align}
\omega_1^2 = -\omega_2^1 = p\omega^1 + q\omega^2, \label{eq:conn_form_def}
\end{align}
由结构方程可解得
$$p = \frac{\mathrm{d}\omega^1}{\omega^1\land\omega^2}=0,\quad q = \frac{\mathrm{d}\omega^2}{\omega^1\land\omega^2}=\frac{1}{u\sqrt{1+u^2}},$$
因此
\begin{align*}
\omega_1^2 = -\omega_2^1 = \frac{1}{u\sqrt{1+u^2}}\omega^2 = \frac{1}{\sqrt{1+u^2}}\mathrm{d}v.
\end{align*}
设法向相关相对分量
$$\omega_1^3 = a\omega^1 + b\omega^2,\quad \omega_2^3 = b\omega^1 + c\omega^2,$$
代入第二基本形式$\mathrm{II}=\omega^1\omega_1^3+\omega^2\omega_2^3$，展开得
\begin{align*}
\mathrm{II}=a(1+u^2)(\mathrm{d}u)^2 + 2bu\sqrt{1+u^2}\mathrm{d}u\mathrm{d}v + cu^2(\mathrm{d}v)^2.
\end{align*}
比较系数：
$$a(1+u^2)=\frac{1}{\sqrt{1+u^2}},\quad 2bu=0,\quad cu^2=\frac{u^2}{\sqrt{1+u^2}},$$
解得
$$a=\frac{1}{(\sqrt{1+u^2})^3},\quad b=0,\quad c=\frac{1}{\sqrt{1+u^2}}.$$
于是
\begin{align*}
\omega_1^3 = -\omega_3^1 = \frac{1}{1+u^2}\mathrm{d}u, \quad
\omega_2^3 = -\omega_3^2 = \frac{u}{\sqrt{1+u^2}}\mathrm{d}v.
\end{align*}
计算外微分：
\begin{align*}
\mathrm{d}\omega_1^2 = -\frac{u}{(\sqrt{1+u^2})^3}\mathrm{d}u\land\mathrm{d}v, \quad
\mathrm{d}\omega_1^3 = 0, \quad
\mathrm{d}\omega_2^3 = \frac{1}{(\sqrt{1+u^2})^3}\mathrm{d}u\land\mathrm{d}v.
\end{align*}
验证结构方程（Gauss-Codazzi方程）：
\begin{align*}
\omega_1^3\land\omega_2^3 &= -\frac{u}{(\sqrt{1+u^2})^3}\mathrm{d}u\land\mathrm{d}v = \mathrm{d}\omega_1^2, \\
\omega_1^2\land\omega_1^3 &= 0 = \mathrm{d}\omega_1^3, \\
\omega_2^2\land\omega_1^3 &= \frac{1}{(\sqrt{1+u^2})^3}\mathrm{d}u\land\mathrm{d}v = \mathrm{d}\omega_2^3.
\end{align*}
即$\mathrm{d}\omega_1^2=\omega_1^3\land\omega_2^3$，$\mathrm{d}\omega_1^3=\omega_1^2\land\omega_2^3$，$\mathrm{d}\omega_2^3=\omega_2^1\land\omega_1^3$，Gauss-Codazzi方程成立.

\end{solution}
\end{example}

\begin{definition}
曲面内蕴微分几何中，一次微分形式$\boldsymbol{\omega_1^2=-\omega_2^1}$称为\textbf{联络形式}.
\end{definition}

\begin{definition}
设$\{\boldsymbol{r};\boldsymbol{e}_1,\boldsymbol{e}_2,\boldsymbol{e}_3\}$为曲面$S$的一阶标架场，满足
$$\mathrm{d}\boldsymbol{e}_1=\omega_1^2\boldsymbol{e}_2+\omega_1^3\boldsymbol{e}_3,\quad \mathrm{d}\boldsymbol{e}_2=\omega_2^1\boldsymbol{e}_1+\omega_2^3\boldsymbol{e}_3.$$
切向量场$\boldsymbol{X}=x^1\boldsymbol{e}_1+x^2\boldsymbol{e}_2$的\textbf{协变微分}定义为
\begin{align*}
\mathrm{D}\boldsymbol{X} &= \mathrm{d}x^1\boldsymbol{e}_1 + x^1\mathrm{D}\boldsymbol{e}_1 + \mathrm{d}x^2\boldsymbol{e}_2 + x^2\mathrm{D}\boldsymbol{e}_2 \\
&= (\mathrm{d}x^1 + x^2\omega_2^1)\boldsymbol{e}_1 + (\mathrm{d}x^2 + x^1\omega_1^2)\boldsymbol{e}_2,
\end{align*}
其中分量的协变微分为
$$\mathrm{D}x^1 = \mathrm{d}x^1 + x^2\omega_2^1,\quad \mathrm{D}x^2 = \mathrm{d}x^2 + x^1\omega_1^2.$$
\end{definition}

\begin{remark}
协变微分仅依赖联络形式$\omega_1^2=-\omega_2^1$，由第一基本形式唯一确定，与第二基本形式无关，是曲面的内蕴几何概念.
\end{remark}

曲面一阶标架场容许绕法向量$\boldsymbol{e}_3$的旋转变换：
\begin{align}
\begin{cases}
\tilde{\boldsymbol{e}}_1 = \cos\theta\boldsymbol{e}_1 + \sin\theta\boldsymbol{e}_2, \\
\tilde{\boldsymbol{e}}_2 = -\sin\theta\boldsymbol{e}_1 + \cos\theta\boldsymbol{e}_2, \\
\tilde{\boldsymbol{e}}_3 = \boldsymbol{e}_3,
\end{cases} \label{eq:frame_trans}
\end{align}
其中$\theta$为曲面上光滑函数.

\begin{theorem}
若一阶标架场经受变换\eqref{eq:frame_trans}，则其相对分量满足
\begin{align}
\begin{cases}
(\tilde{\omega}^1,\tilde{\omega}^2) = (\omega^1,\omega^2)\begin{pmatrix}
\cos\theta & -\sin\theta \\
\sin\theta & \cos\theta
\end{pmatrix}, \\
\tilde{\omega}^3 = \omega^3 = 0, \\
\tilde{\omega}_1^2 = \omega_1^2 + \mathrm{d}\theta, \\
(\tilde{\omega}_1^3,\tilde{\omega}_2^3) = (\omega_1^3,\omega_2^3)\begin{pmatrix}
\cos\theta & -\sin\theta \\
\sin\theta & \cos\theta
\end{pmatrix}.
\end{cases} \label{eq:comp_trans}
\end{align}
\end{theorem}

\begin{proof}

由$\mathrm{d}\boldsymbol{r}=(\omega^1,\omega^2)\begin{pmatrix}\boldsymbol{e}_1\\\boldsymbol{e}_2\end{pmatrix}=(\tilde{\omega}^1,\tilde{\omega}^2)\begin{pmatrix}\tilde{\boldsymbol{e}}_1\\\tilde{\boldsymbol{e}}_2\end{pmatrix}$，代入标架变换式可得$(\tilde{\omega}^1,\tilde{\omega}^2)$的变换规律.
由$\tilde{\boldsymbol{e}}_3=\boldsymbol{e}_3$，得$\tilde{\omega}^3=\omega^3=0$.
联络形式：
\begin{align*}
\tilde{\omega}_1^2 &= \mathrm{D}\tilde{\boldsymbol{e}}_1\cdot\tilde{\boldsymbol{e}}_2 \\
&= \mathrm{D}(\cos\theta\boldsymbol{e}_1+\sin\theta\boldsymbol{e}_2)\cdot(-\sin\theta\boldsymbol{e}_1+\cos\theta\boldsymbol{e}_2) \\
&= \mathrm{d}\theta + \omega_1^2.
\end{align*}
法向分量变换由$\mathrm{d}\tilde{\boldsymbol{e}}_3=\mathrm{d}\boldsymbol{e}_3$直接可得.

\end{proof}

\begin{example}
利用一阶标架场的相对分量，构造定义在曲面$S$上、与标架选取无关的二次微分式.
\begin{solution}

由定理\eqref{eq:comp_trans}，二次微分式
$$(\omega^1)^2+(\omega^2)^2,\quad \omega^1\omega_1^3+\omega^2\omega_2^3,\quad (\omega_1^3)^2+(\omega_2^3)^2$$
均与标架选取无关，依次对应曲面的第一、第二、第三基本形式.

\end{solution}
\end{example}

\begin{example}
证明二次外微分式$\omega^1\land\omega^2$与一阶标架场的选取无关.
\begin{proof}

\begin{align*}
\tilde{\omega}^1\land\tilde{\omega}^2 &= (\cos\theta\omega^1+\sin\theta\omega^2)\land(-\sin\theta\omega^1+\cos\theta\omega^2) \\
&= \omega^1\land\omega^2.
\end{align*}
$\omega^1\land\omega^2=\sqrt{EG-F^2}\mathrm{d}u\land\mathrm{d}v$，即为曲面的面积元素$\mathrm{d}\sigma$，与标架选取无关.

\end{proof}
\end{example}

\begin{example}
证明$\mathrm{d}\omega_1^2$，$\omega_1^3\land\omega_2^3$，$\omega^1\land\omega_2^3-\omega^2\land\omega_1^3$是与标架选取无关的二次外微分式.
\begin{proof}

由定理\eqref{eq:comp_trans}，$\tilde{\omega}_1^2=\omega_1^2+\mathrm{d}\theta$，故
$$\mathrm{d}\tilde{\omega}_1^2=\mathrm{d}\omega_1^2+\mathrm{d}(\mathrm{d}\theta)=\mathrm{d}\omega_1^2.$$
对于法向分量：
\begin{align*}
&\tilde{\omega}_1^3\land\tilde{\omega}_2^3 = (\cos\theta\omega_1^3+\sin\theta\omega_2^3)\land(-\sin\theta\omega_1^3+\cos\theta\omega_2^3) = \omega_1^3\land\omega_2^3, \\
&\tilde{\omega}^1\land\tilde{\omega}_2^3-\tilde{\omega}^2\land\tilde{\omega}_1^3 = \omega^1\land\omega_2^3-\omega^2\land\omega_1^3.
\end{align*}
设$\omega_1^3=a\omega^1+b\omega^2$，$\omega_2^3=b\omega^1+c\omega^2$，则
\begin{align*}
&\omega_1^3\land\omega_2^3 = (ac-b^2)\omega^1\land\omega^2, \\
&\omega^1\land\omega_2^3-\omega^2\land\omega_1^3 = (a+c)\omega^1\land\omega^2.
\end{align*}
因$\omega^1\land\omega^2$为不变量，故$ac-b^2$、$a+c$为曲面不变量.

\end{proof}
\end{example}

\begin{definition}
曲面的\textbf{Gauss曲率}$K$与\textbf{平均曲率}$H$定义为
$$K=ac-b^2=\frac{LN-M^2}{EG-F^2},\quad H=\frac{1}{2}(a+c)=\frac{1}{2}\left(\frac{L}{E}+\frac{N}{G}\right).$$
\end{definition}

\begin{theorem}
Gauss曲率$K$是曲面的内蕴不变量，仅由第一基本形式决定，表达式为
\begin{align}
K = -\frac{\mathrm{d}\omega_1^2}{\omega^1\land\omega^2}. \label{eq:gauss_curv}
\end{align}
\end{theorem}

\begin{proof}

由Gauss结构方程$\mathrm{d}\omega_1^2=\omega_1^3\land\omega_2^3=-K\omega^1\land\omega^2$，变形即得\eqref{eq:gauss_curv}.
因$\omega_1^2,\omega^1,\omega^2$均由第一基本形式唯一确定，故$K$为内蕴不变量.

\end{proof}

\begin{example}
设曲面$S$的第一基本形式为$\mathrm{I}=E(\mathrm{d}u)^2+G(\mathrm{d}v)^2$，证明曲面$S$的Gauss曲率满足
\begin{align}
K=-\frac{1}{\sqrt{EG}}\left(\left(\frac{(\sqrt{E})_v}{\sqrt{G}}\right)_v+\left(\frac{(\sqrt{G})_u}{\sqrt{E}}\right)_u\right).
\label{eq:gauss-curvature-orth}
\end{align}
\end{example}

\begin{proof}

将曲面$S$的第一基本形式配平方为
\begin{align*}
\mathrm{I}=(\sqrt{E}\mathrm{d}u)^2+(\sqrt{G}\mathrm{d}v)^2,
\end{align*}
取曲面$S$上一阶标架场的相对分量
\begin{align*}
\omega^1=\sqrt{E}\mathrm{d}u,\quad \omega^2=\sqrt{G}\mathrm{d}v.
\end{align*}
由前文结论可得
\begin{align*}
\omega_1^2=-\frac{(\sqrt{E})_v}{\sqrt{G}}\mathrm{d}u+\frac{(\sqrt{G})_u}{\sqrt{E}}\mathrm{d}v.
\end{align*}
直接计算外微分与外积：
\begin{align*}
\mathrm{d}\omega_1^2=\left(\left(\frac{(\sqrt{E})_v}{\sqrt{G}}\right)_v+\left(\frac{(\sqrt{G})_u}{\sqrt{E}}\right)_u\right)\mathrm{d}u\wedge\mathrm{d}v,\\
\omega^1\wedge\omega^2=\sqrt{EG}\mathrm{d}u\wedge\mathrm{d}v.
\end{align*}
结合Gauss曲率定义$K=-\frac{\mathrm{d}\omega_1^2}{\omega^1\wedge\omega^2}$，代入即得式\eqref{eq:gauss-curvature-orth}。

\end{proof}

\begin{example}
证明二次微分形式$\omega^1\omega_2^3-\omega^2\omega_1^3$与曲面$S$上一阶标架场的选取无关。
\end{example}

\begin{proof}

一阶标架场的正交变换满足
\begin{align*}
(\widetilde{\omega}^1,\widetilde{\omega}^2)=(\omega^1,\omega^2)\cdot\begin{pmatrix}\cos\theta&-\sin\theta\\\sin\theta&\cos\theta\end{pmatrix},\\
\begin{pmatrix}\widetilde{\omega}_2^3\\-\widetilde{\omega}_1^3\end{pmatrix}=\begin{pmatrix}\cos\theta&\sin\theta\\-\sin\theta&\cos\theta\end{pmatrix}\begin{pmatrix}\omega_2^3\\-\omega_1^3\end{pmatrix}.
\end{align*}
因此
\begin{align*}
(\widetilde{\omega}^1,\widetilde{\omega}^2)\begin{pmatrix}\widetilde{\omega}_2^3\\\widetilde{\omega}_1^3\end{pmatrix}=(\omega^1,\omega^2)\begin{pmatrix}\omega_2^3\\-\omega_1^3\end{pmatrix}.
\end{align*}
将$\omega_1^3=a\omega^1+b\omega^2,\ \omega_2^3=b\omega^1+c\omega^2$代入上式，得
\begin{align}
\omega^1\omega_2^3-\omega^2\omega_1^3=b(\omega^1)^2+(c-a)\omega^1\omega^2-b(\omega^2)^2.
\label{eq:invariant-2form}
\end{align}
该表达式在标架正交变换下形式不变，故该二次微分形式与一阶标架场的选取无关。

\end{proof}

\begin{example}
设$\{\bm{r};\bm{r}_u,\bm{r}_v,\bm{n}\}$是曲面$S$上的自然标架场，其相对分量$\omega^j,\omega_i^j$如(4.3)和(4.4)式给出。证明：
\begin{enumerate}[(1)]
\item $\mathrm{d}g_{\alpha\beta}=g_{\alpha\gamma}\omega_\beta^\gamma+g_{\gamma\beta}\omega_\alpha^\gamma$;
\item 命$\omega_{\alpha\beta}=\omega_\alpha^\gamma g_{\gamma\beta}$，$R_{\alpha\beta\gamma\delta}=g_{\beta\eta}R_{\alpha\gamma\delta}^\eta$，则
\begin{align*}
\mathrm{d}\omega_{\alpha\beta}+\omega_\alpha^\gamma\land\omega_{\beta\gamma}=-\frac{1}{2}R_{\alpha\beta\gamma\delta}\mathrm{d}u^\gamma\land\mathrm{d}u^\delta;
\end{align*}
\item $\mathrm{d}g^{\alpha\beta}=-g^{\alpha\gamma}\omega_\gamma^\beta-g^{\gamma\beta}\omega_\gamma^\alpha$.
\end{enumerate}
\end{example}

\begin{example}
设曲面的第一基本形式给定如下：
\begin{enumerate}[(1)]
\item $\mathrm{I}=\frac{1}{v^2}\left((\mathrm{d}u)^2+(\mathrm{d}v)^2\right),\ v>0$;
\item $\mathrm{I}=\frac{(\mathrm{d}u)^2-4v\mathrm{d}u\mathrm{d}v+4u(\mathrm{d}v)^2}{4\left(u-v^2\right)},\ u>v^2$;
\item $\mathrm{I}=(\mathrm{d}u)^2+2\cos\psi\mathrm{d}u\mathrm{d}v+(\mathrm{d}v)^2$，$\psi$是$u,v$的连续可微函数.
\end{enumerate}
求该曲面上关于一个一阶标架场的相对分量$\omega^1,\omega^2,\omega_1^2$和Gauss曲率$K$.
\end{example}

\begin{example}
在旋转曲面
\begin{align*}
\bm{r}(u,v)=(u\cos v,u\sin v,g(u))
\end{align*}
上建立一阶标架场，并计算它的相对分量$\omega^j,\omega_i^j$.
\end{example}

\begin{example}
对于例题2所给出的一阶标架场的相对分量$\omega^j,\omega_i^j$，写出它们所满足的第二组结构方程.
\end{example}

\begin{example}
已知曲面$S$的两个基本形式分别为
\begin{align*}
\mathrm{I}=(a^2+2v^2)(\mathrm{d}u)^2+4uv\mathrm{d}u\mathrm{d}v+(a^2+2u^2)(\mathrm{d}v)^2,
\end{align*}
\begin{align*}
\mathrm{II}=-\frac{2a}{\sqrt{a^2+2\left(u^2+v^2\right)}}\mathrm{d}u\mathrm{d}v.
\end{align*}
求该曲面上一个一阶标架场的相对分量$\omega^j,\omega_i^j$，并且验证它们适合第二组结构方程.
\end{example}



\end{document}